%! AP Statistics — Unit 1, Lesson 1.3 — Group Activity ANSWER KEY
\documentclass[10pt]{article}
\usepackage{apstats-article}
\usepackage{apstats-key}

%=============================================================================
\begin{document}

\pageheader{Unit 1, Lesson 1.3}{Group Activity --- Answer Key}
\begin{tabularx}{\linewidth}{X X X}
Name:\;\ans{Key} & Partner:\; & Period:\;
\end{tabularx}

\vspace{0.1in}

% ── PART 1 ───────────────────────────────────────────────────────────────────
{\large\bfseries\color{navy} Part 1 \quad Build the Table}

\vspace{0.1in}

\small\noindent Raw data: 30 middle school students, favorite sport.\\
Counts: Baseball~5, Basketball~7, Soccer~9, Swimming~4, Tennis~5.

\vspace{0.1in}

\renewcommand{\arraystretch}{2.0}
\begin{tabularx}{\linewidth}{@{} >{\bfseries}p{0.22\linewidth}
                                    C{0.18\linewidth}
                                    C{0.2\linewidth}
                                    C{0.24\linewidth} @{}}
\rowcolor{sky}
\textbf{Sport} & \textbf{Tally} & \textbf{Frequency} & \textbf{Rel.\ Frequency} \\
\midrule
Baseball   & \small{|||||} & \ans{5} & \ans{5/30 $\approx$ 0.167 (16.7\%)} \\
Basketball & \small{|||||||} & \ans{7} & \ans{7/30 $\approx$ 0.233 (23.3\%)} \\
Soccer     & \small{|||||||||} & \ans{9} & \ans{9/30 = 0.300 (30.0\%)} \\
Swimming   & \small{||||} & \ans{4} & \ans{4/30 $\approx$ 0.133 (13.3\%)} \\
Tennis     & \small{|||||} & \ans{5} & \ans{5/30 $\approx$ 0.167 (16.7\%)} \\
\midrule
\rowcolor{sky}\textbf{Total} & & \ans{30} & \ans{1.000} \\
\end{tabularx}

\begin{teachernote}
Verify students' raw data tallies match the given list. The data in the activity
was designed so the counts above are correct. Accept rounding to 3 decimal places.
\end{teachernote}

\vspace{0.1in}

\begin{enumerate}[leftmargin=1.5em, itemsep=8pt]
  \item Total = 30? \ans{Yes.}\quad Sum $= 1$? \ans{Yes (allow small rounding).}

  \item Most popular: \ans{Soccer, chosen by 9 students (30\%).}

  \item \ans{Of the 30 middle school students surveyed, soccer was the most popular
        sport (30\%). Swimming was the least popular (13.3\%), with baseball and
        tennis tied as the second least popular (16.7\% each).}
\end{enumerate}

\vspace{0.12in}

% ── PART 2 ───────────────────────────────────────────────────────────────────
{\large\bfseries\color{navy} Part 2 \quad Compare Two Groups}

\vspace{0.1in}

\renewcommand{\arraystretch}{2.0}
\begin{tabularx}{\linewidth}{@{} >{\bfseries}p{0.22\linewidth}
                                    C{0.18\linewidth}
                                    C{0.2\linewidth}
                                    C{0.24\linewidth} @{}}
\rowcolor{sky}
\textbf{Sport} & \textbf{Frequency} & \multicolumn{2}{c}{\textbf{Rel.\ Frequency}} \\
\midrule
Baseball   & 8  & \multicolumn{2}{c}{\ans{0.160 (16.0\%)}} \\
Basketball & 15 & \multicolumn{2}{c}{\ans{0.300 (30.0\%)}} \\
Soccer     & 14 & \multicolumn{2}{c}{\ans{0.280 (28.0\%)}} \\
Swimming   & 7  & \multicolumn{2}{c}{\ans{0.140 (14.0\%)}} \\
Tennis     & 6  & \multicolumn{2}{c}{\ans{0.120 (12.0\%)}} \\
\midrule
\rowcolor{sky}\textbf{Total} & 50 & \multicolumn{2}{c}{\ans{1.000}} \\
\end{tabularx}

\vspace{0.1in}

\begin{enumerate}[leftmargin=1.5em, itemsep=8pt]

  \item \ans{Not fully valid. Soccer does appear near the top in both groups, but
        the group sizes differ (30 vs.\ 50), so we need relative frequencies to
        compare fairly. In the middle school group, 30\% chose soccer; in the high
        school group, 28\% chose soccer --- the difference is small.}

  \item \ans{High school shows a stronger preference for basketball: 30\% (15/50)
        vs.\ 23.3\% (7/30) in middle school.}

  \item \ans{Third school: $60/200 = 0.30 = 30\%$. The middle school group also
        has 30\% choosing soccer, so soccer is equally popular in both groups by
        relative frequency, even though the third school has more students.}

\end{enumerate}

\vspace{0.12in}

% ── PART 3 ───────────────────────────────────────────────────────────────────
{\large\bfseries\color{navy} Part 3 \quad Reflection}

\vspace{0.1in}

\begin{tcolorbox}[colback=goldbg, colframe=goldacc, boxrule=0.8pt, arc=2mm,
    left=3mm, right=3mm, top=2mm, bottom=2mm]
\small
\ans{A relative frequency table expresses each count as a proportion of the total,
which makes comparisons fair regardless of group size. For example, comparing
``9 soccer players'' (middle school, $n=30$) to ``14 soccer players''
(high school, $n=50$) is misleading --- high school has more simply because it
is a larger group. Using relative frequencies (30\% vs.\ 28\%) shows that
soccer is roughly equally popular in both groups.}
\end{tcolorbox}

\end{document}
